\documentclass[11pt]{article}
\usepackage[utf8]{inputenc}
\usepackage[T1]{fontenc}
\usepackage{booktabs}
\usepackage{amsmath}
\usepackage{geometry}
\geometry{margin=2.5cm}

\title{Causal Fairness Analysis Report}
\author{Generated by FairMind and an LLM}
\date{2026-08-22}

\begin{document}
\maketitle

\section{Analysis setup}

\begin{tabular}{ll}
\toprule
Dataset & oulad\_final\_result \\
Protected attribute ($X$) & S1\_gender \\
Comparison & F $\to$ M \\
Outcome ($Y$) & T\_final\_result (Pass) \\
Mediator ($W$) & studied\_credits \\
Confounder ($Z$) & highest\_education \\
\bottomrule
\end{tabular}

\section{Causal effects (FairMind, exact values)}

The five effects below are computed by FairMind through exact inference on the
fitted Bayesian network. These values are final and must not be recomputed or
altered.

\begin{tabular}{lr}
\toprule
Effect & Value \\
\midrule
Total Variation (TV) & -0.027884 \\
Total Effect (TE) & -0.036580 \\
Spurious Effect (SE) & 0.008696 \\
Direct Effect (DE) & -0.034469 \\
Indirect Effect (IE, decomposition form: TE = DE + IE) & -0.002111 \\
\bottomrule
\end{tabular}

\section{Qualitative interpretation}

\subsection{Total effect and spurious component (TV, TE, SE)}

The total disparity between the two groups is relatively small, with a total variation of -0.027884. After removing the confounding effect of the highest education level, the genuine causal effect (TE) is -0.036580, indicating a significant portion of the observed disparity is due to genuine causal factors. However, a portion of the observed disparity is spurious, driven by the confounder highest education level (SE = 0.008696), which is relatively small compared to the total variation.

\subsection{Direct effect (DE)}

The direct effect (DE) is -0.034469, indicating a strong direct discrimination against the protected group (F to M). This direct effect is substantial and suggests that the protected group faces a significant disadvantage that is not mediated through the studied credits.

\subsection{Indirect effect (IE)}

The indirect effect (IE) is -0.002111, which is relatively small compared to the direct effect. Since the indirect effect has the same sign as the direct effect, it amplifies the total effect, making the overall disadvantage even more pronounced. The mediator channel (studied credits) does not mitigate the direct effect but rather adds to it.

\section{Recap Questions}

\begin{enumerate}
\item Is there direct discrimination against the protected group? \textbf{Answer:} YES
\item Does the indirect effect through the mediator mitigate the total effect? \textbf{Answer:} NO
\item Does the observed total effect exceed the practical relevance threshold? \textbf{Answer:} NO
\item Does the spurious component (SE) contribute substantially to the observed difference? \textbf{Answer:} NO
\item Is the direct effect the dominant component compared to the indirect effect? \textbf{Answer:} YES
\end{enumerate}

\vspace{1em}
\noindent\footnotesize
The recap answers follow deterministic threshold rules applied to the values above: direct discrimination when $|\mathrm{DE}| > 0.05$; a mediator channel counted as negligible when $|\mathrm{IE}| < 0.005$; practical relevance when $|\mathrm{TV}| > 0.05$; a substantial spurious component when $|\mathrm{SE}| / |\mathrm{TV}| > 0.1$.

\end{document}